\documentclass{build/zhvt-classic}
\title{國劇身段譜}[國~劇~身~段~譜]
\author{齊如山}
\maker{甲辰年【龙】腊月十三酉时【日入】重製}

\begin{document}

\maketitle


\grid{none}{none}

\insertgraphic[width=0.65\linewidth,angle=90]{pictures/ti1}

\insertgraphic[width=0.65\linewidth,angle=90]{pictures/ti2}

\insertgraphic[width=0.65\linewidth,angle=90]{pictures/ti3}

\insertgraphic[width=0.65\linewidth,angle=90]{pictures/ti4}

\insertgraphic[width=0.65\linewidth,angle=90]{pictures/ti5}

\insertgraphic[width=0.65\linewidth,angle=90]{pictures/ti6}

\insertgraphic[width=0.65\linewidth,angle=90]{pictures/ti7}

\insertgraphic[width=0.65\linewidth,angle=90]{pictures/ti8}

\insertgraphic[width=0.65\linewidth,angle=90]{pictures/ti9}

\insertgraphic[width=0.65\linewidth,angle=90]{pictures/ti10}

\clearpage
\gridall

\tableofcontents
%  序
\input{chapters/chap-0.tex}

%  第一章 論剧來源於古之歌舞
\input{chapters/chap-1.tex}
%  第二章 論放劇與唐朝之舞有密切關係
\chapter*[第二章]{論放劇與唐朝之舞有密切關係}{}   

现在戏剧的身段、固然是由古代之舞嬗变而来。但是由唐朝起、才有了具体的规模。

涵虚子论剧云
\begin{preface}

﹁戏剧至隋朝始盛、隋谓之康衢戏、唐谓之梨园乐、宋谓之华林戏、元谓之升平乐﹂云云。

\end{preface}

陶九成论曲云

\begin{preface}

﹁唐有传奇、宋有戏曲、金有院本雑剧、而元因之。然院制杂剧、厘而为二矣﹂云云。

\end{preface}

以上二说、一论戏剧、一论词曲、意思大致相同。涵虚子虽然说是戏剧至隋朝始盛、然隋朝为期甚短、不见得有完备的组织、当然是到了唐朝、经有大力量的唐明皇一提倡、才有了完备的规定。以后由宋而元而明清、虽一朝与一朝的名目不同、但是没有什么极大的变更一直传流到了现在、各朝之命名、虽然屡经更改。【清明管戏剧的衙门亦名曰昇平*】而戏界至今、仍日【梨园行】就是戏中的各种身段步法、大致也来源于唐朝之舞。按古人歌舞曲牌的名子、存留者尚不少、而怎样的舞法、各种记载中、则毫未提及、无从参考。惟有宋朝周密的癸辛杂识里头、有一段记载说是﹁予尝得故都德寿宫舞谱二大帙、其中皆新制曲、多妃嫔诸阁、分所进者。所谓讲者、其间在所谓：

~左右垂手\quad 雙拂\quad 抱肘\quad 合蝉\quad 小转\quad 虚影\quad 横影\quad 稱裹 

~大小转攛\quad 盤转\quad 叉腰\quad 捧心\quad 叉手\quad 打场\quad 搀手\quad 鼓兒

~打鸳鸯场\quad 分颈\quad 回头\quad 海眼\quad 收尾\quad 豁头\quad 舒手\quad 布过

~鲍老惙\quad  对窠\quad  方腾\quad  齐收\quad  舞头\quad  舞尾\quad  呈手\quad  关宝


~掉袖儿\quad  拂\quad  躜\quad  绰\quad  觑\quad  掇\quad  蹬\quad  焌


~五花儿\quad  踢\quad  搕\quad  剌\quad  𭣇\quad  繫\quad  搠\quad  捽


~雁翅儿\quad  靠\quad  挨\quad  拽\quad  捺\quad  闪\quad  缠\quad  提


~龟背儿\quad  踏\quad  儹\quad  木\quad  摺\quad  促\quad  当\quad  前


~勤蹄儿\quad  摆\quad  磨\quad  捧\quad  抛\quad  奔\quad  抬\quad  撅

以上唐代之舞谱也。其中所云之左右垂手、大小转攛、打鸳鸯场等九种名目。乃是舞牌之名。至雙拂、盤转、分颈、对窠等等六十三种名目、乃各种姿式之名词也。只看他将各种舞的姿式、又特定名目、便可知彼时歌舞之盛行、及衆【众】人研究之状况。当年此种舞谱、一定很多、可惜存留到今者、只有一二页、可谓吉光片羽矣。但是我们拿这几十个名词、与现在戏中之身段来参证、就觉着已经有许多相合的地方。是可知戏剧之身段、实来源于唐代之舞也。现在把谱中姿式的名词、与戏中各种身段姿式来参证参证：

\quad 左右垂手\quad 

这个舞牌中姿式的名词、于现在戏中旦脚出场的各种身段、大致相似且读古来各种记载和歌咏的文字、大致都说这个舞牌子是女子所舞的；则旦脚出场之身段、来源于此也未可知。因旦脚所去的都是女子也。

兹参证如下比方：

\quad 双拂\quad  即戏中之双抖袖、因拂字之义、舆抖字大致相同也。

\quad 抱肘\quad  即戏中之合抱、俗名抱夹、乃双手交于胸际。

\quad 合蝉\quad  即戏中之双背袖即双贫手、刚袖往后一背、如同蝉之合翼。

\quad 小转 \quad
即戏中之半转身、即往后看。按宋朝朱熹诗傅释﹂辗转反侧四字之意、曰辗者转之、转者辗之周、反者转之过、侧者转之留﹂舞谱亦用此字样、日转初势、转半势、转周势、转过势、转留势、伏覩势、仰瞻势、回头势。摢此则小转之意、当然要外乎回头、和半转身的情形。  

\quad 虚影
\quad 
即戏中之掍、或一掍、或两掍、即然曰虚影、当然不是呆板的姿式、所以说他是戏中之掍。

\quad 横影
\quad 即戏中之又舒袖、即两手旁伸。從前旦脚此種身段颇多、梆子班中、花旦尤多。又手往旁一伸、则横*于高、与横影二字的意思𣑕合。

\quad 稱裹
\quad 即戏中之左右理装、俗名左右看、当看左旁的时候、则將右袖举起、用左袖作担鞋现形式。看右旁的时候反是、所谓稱裹者、大致係看看装裹的相稱、不相稱的意思。不则裹字在姿式名词中、用之殊觉生硬。

\quad 掉䄂儿
\quad 这个舞的牌子、按古人文字中考之、也是女子所舞的、且与旦脚出场之各種身段、也大致相似。比方：

\quad 拂
\quad 即戏中之單【单】抖袖。

\quad 躦
\quad 即戏中之蹲、類篇云﹁𧿙、聚足、或作蹲。﹂又云﹁踏也。﹂戏中每合抱、必蹲身子、證以類篇所谓聚足又踏也等语、或亦係旦脚之踏步。即左足斜置于右足之后、而两腿聚靠、或右足斜置于左足之后、旦脚亮像必须如此。

\quad 綽
\quad 即戏中之欲前走后退一二步之身段、所谓绰約也。

\quad 覷
\quad 即旦脚之用手按鬓角、而往前看之姿式。旦角出场立定时、必有这种身段。

\quad 掇\quad 
即戏中之提鞋式、或蹲身大抖袖式、此为旦脚出场必有的身段。用右手在右腿后提左鞋、用左手在左腿后提右鞋、花旦则用手掇其鞋后根、青衣则略以见意而已。

\quad 蹬\quad 
即戏中之超步、所谓𨃴蹬也。

\quad 焌\quad 
疑即戏中之双举袖前行式、旦脚行路、必打腰裙、在匆或着急的时候、必将袖双举。比方牧羊圈讨饭一场、青衣与老旦出场时、青衣之身段、便用此式。所谓燃烧火*上之炎状也。又周禮注﹁杜子春曰：焌與俊通。﹂则亦是正面与人看之义。此二釋義、与身段都算相合。

\quad 大小转撺\quad 


这个舞牌、在记载中不多见、揣摩他舞的情形、似一种群舞。但是各人的身段、都是相同的。且其名词之情形、亦与日脚之身段、有相似的地方。但此与前左右垂手、掉袖儿两种有不同之处。盖前两种之姿式、差不多与旦脚出场之情形、完全相同、可以说旦脚出场之姿式、完全是采用的那两种舞、至于此种舞牌、虽与旦脚之身段、有相合之处；但系散见于各戏之中、不似前两种之整个的采用也。现在也把他来参证参证：

\quad 盘转
\quad 疑即戏中之绕圆场、现在仍名日走圆场、或跑圆场。如穆柯寨穆桂英登高台之前、及射雁后进场之前、都有此种身段。且绕圆场三字的意思、与盘转二字也相同。

\quad 叉腰
\quad 
即戏中之叉腰、双手叉于腰间、手背朝裹、肘往前伸、此种身段、旦脚用之、非常美观。

\quad 捧心
\quad 
即双手重叠于胸前、现在仍名日捧心、但亦旦脚用之。

\quad 叉手
\quad 
此名词颇难解释、疑系多人互相叉、手戏中不多见。如以个人叉手解释、则係双手互叉、手心向下之状。旦脚往往有此身段、花旦尤多、但此不过在疑似之间。

\quad 打场 
\quad 
疑即戏中之旋转、因*家打场、皆係牵引碜碡旋转也、河北童谣、有打场谣曰﹁咧咧咧咧场、打喽𢾶君打高梁。﹂等语。小儿每逢念此、身子必旋转、而每旋转、亦必歌也。此谣行之颇广、来源当亦甚远。故疑打场二字、即旋转之义。戏中旋转之处甚多、亦有两人互转之时。二人各以左膊相跨而转、如金山寺青蛇白蛇、及牧羊圈青衣老旦讨饭一场、都有这种身段。

\quad 搀手 
\quad 
疑即戏中二人两手互搀、如金山寺青蛇白蛇、及奇双会赵宠合李桂枝等等、都有此身段。惟此种身段、有相哭时方用之。

\quad 鼓儿
\quad 
阙疑

\quad 勤蹄儿
\quad 
这个舞牌、虽然以蹄字为名、但谱中名词、除奔字身外、大致都是用手的名词、所以鄙人疑惑此种舞系用袖子之舞、因谱中名词与现在戏中用袖子之之词大致相同也。

比方：

\quad 摆
\quad 
即戏中之摆袖、戏中的名词、用袖先往裹一转、再往外摆者名抖袖。一直往外一摆者、日名摆袖。

\quad 磨
\quad 
即戏中之磨袖、似担靴的样而确非担靴、比方以左袖平举而拍右腿、即以右袖绕转于腿之旁。如张飞及火判、常有这种身段。

\quad 捧 
\quad 
即戏中之捧袖、如双袖先往裹一折、作拱手状、此即名日捧袖。

\quad 抛
\quad 
即戏中之抛袖。抛袖之形势、不止一种、如双袖同时往前一挥、或两袖错棕挥之。总之戏中之抛袖、投袖、挥袖、掷袖等等、都是抛的性质、此说参看第三章。

\quad 奔 
\quad 
此字颇难解释、或係往前抛袖取奔放之义、因若以字面解释则只有快跑一解。在一小小舞场之中、也似乎用不着快跑的身段。

\quad 擡
\quad 
即戏中之撩袖、于惊讶、恐惧、大哭时、往往用此身段。其式亦不一、详见第三章。

\quad 撅
\quad 
即戏中拈袖之意、戏中有抓袖、拈袖、折袖、攘袖等、皆系撅袖的意思。详见第三章。

\quad 龟背儿
\quad 
这个舞牌、在记载中、亦不多见、不知是什么意思、按唐朝段安节乐府雜录中、有：字舞、花舞等名目系用舞的人、排成一个字、或排成一朵花的意思。如此则此舞牌、或者是排成一个龟背的形式、也未可知。但看此谱中的名词、亦系多数人的舞法、好像与现在戏中堆花、堆鬼、摆莲花灯、摆七巧阃等等的舞法、大致相同。

\quad 踏
\quad 
疑即戏中之踏步。右足迈于左腿之左、或左足迈于右腿之右、皆名曰踏步。此系生净等男子的身段。旦脚踏步尤其多遇亮像时、必须踏步、就是将左足聚立于右足之后、或右足坚立于左足之后、都是踏步。再者戏中多数人舞时、往往此人跪于地上、又一人用足踏于跪者之腿上、此在群舞中乃极常用之身段。或与此踏字有直接关系也未可知。

\quad 儹
\quad 
正韵云﹁儹者、聚也。﹂疑即戏中之攒字、因儹攒二字、音义都相同也。如今成中、只要大家在一处凑聚、便名曰攒。武行起打、一人挡四人的时候、就叫结攒；俗名又叫结个攒、则攒字完全成一名词矣。与此儹字必有相当的关系也。

\quad 木
\quad 
疑系直立如木的意思。戏中多数人的身段、往往彼此对立、亦名日亮像。

\quad 摺【折】
\quad 
即戏中之翻。如数人舞、一行直走的时候、忽然自第一人起、往回折走、即名日翻；或日回翻、即系折叠的意思。按广韵曰﹁摺亚也。﹂或系叠架之意。戏中也有这样的舞法、如堆灯、堆花等等、往往将切末叠于一处、就是一排比一排高的意思。总之以上两种解释法、一为横叠、一为竖叠、大致都与戏中舞式相合。

\quad 促
\quad 
集韻曰﹁迫也；近也；密也﹂如此解法便是往一处凑集的意思。按戏中群舞、如摆七巧图、或四人武舞等等的时候。有时众人提四门斗、凑在一处、名曰合拢；亦曰凑。若分开、则名曰拉开。按合拢之名词、舆促字的意思颇相合。

\quad 当
\quad 
此字当然是舞者立定、以迎来者的意思、与戏中拉开时之身段颇相合。

\quad 當
\quad 
此字当然是舞者立定、以迎来者的意思、与戏中拉开时之身段颇相合。盖拉开时、必各拉山膀、又武舞。亮像之形势与此亦合。

\quad 前
\quad 
此当然是一齐向前的意思、或者就是戏中之堆花、一排一排的依次向前走的样式。

\quad 雁翅儿
\quad 
舞牌名曰雁翅、想来是所舞的姿式、彷彿大雁的翅膀一样的意思；各种姿式、即像翅膀、则舞者大致不是一个人、其数必在一人以上。且看其舞谱名词、也像两个人、或幾个人合舞的意思、于戏中两人同作的身段、相合的地方很多。而转、姿式亦极美观。或彼此用肩相靠。戏中此种身段亦颇多。

\quad 挨
\quad 
说文﹁挨、击背也。﹂又正字通﹁凡物相近、谓之挨。﹂扬子方言谓﹁强进曰挨。﹂按说文是打的意思、正字通与扬子方言是挤的意思。戏中二人对舞时、用手相击、名曰搕。如金山寺、青蛇白蛇对唱时、恒彼此以剑对搕等等形势合意、都与正字通的解释法相同。愚常疑此种身段之不甚合道理、二人既非互打、又非练习、为什么要彼此相击呢？其来源大致与此有关。又戏中旦脚、用双手叉腰、横行俏步、二人对行、行至一处、两肩相摩、即名曰挨、如佳期红娘唱﹁待红娘轻轻悄悄把门挨﹂时、即系双手叉腰、俏步横行、想这就是形容挨字的意思。这种情形、与正字通及方言之解释法、亦大致相合。

\quad 拽
\quad 
即戏中二人、此人用右手拉彼人之左手、同作卧雲式之身段。因此种身段、实是彼此相拽、且所作身段之姿式、亦如鸟之展翅、与舞牌之名亦相合。

\quad 捺
\quad 
捺者、按也。疑即放中二人之亮高矮像。此人蹲下、彼人立于后边、以手按前人之肩。群舞之时、此种身段尤其多。

\quad 闪
\quad 
即戏中之闪开。比方二人同作身段、走至相近之处、必要往外一躲、又名曰往外一闪。

\quad 缠
\quad 
疑即戏中之二人拉手、同作鹆子翻身、因此即系缠绕的意思。再如二人用枪打仗时、两枪头互搅、即名曰缠、亦名曰绕。

\quad 提
\quad 
戏中蹲身之后、身子起的时候、腿亦随之提起、即名日提。或此人按彼人之肩、身子往上一竄【窜】亦名曰提翻筋斗、有一种名单提。以上三种或与此提字有点关系。

\quad 鲍老惙
\quad 
鲍老惙三字之义意、不大明了、但其中之名词、则完全是多数人合舞之意、与戏中四人合作之身段情形、有相同之点。

\quad 对窠
\quad 
即戏中数人、合凑而各作同鲜身段的意思。

\quad 方腾
\quad 
比名词在戏中、亦颇普通、与四门斗亦大致相同。

\quad 齐收
\quad 
疑即戏中了舞时、各囗本位的意思、戏中即名曰囗本位。按囗本位三字、与齐收二字、亦大致相同。

\quad 舞头
\quad 
阙疑。

\quad 舞尾
\quad 
阙疑。 

\quad 呈手
\quad 
疑即戏中群舞时、各作供手之式。

\quad 关宝
\quad 
此二字、殊无索解。

\quad 五花儿
\quad 
这一种舞牌、一定是各持器械、与戏中之舞如对枪、对刀等场子之姿式、相合。

\quad 踢
\quad 
或即係戏中二人相踢、或一踢一拍。

\quad 榼
\quad 
戏中用刀枪、彼此相撀、名曰搕。

\quad 刺
\quad 
戏中用枪直刺方之人、即名曰刺。

\quad 𭣇
\quad 
戏中甲用枪刀、压住乙之枪刀、乙必用手极力𭣇之、此即名曰𭣇。再一人持枪、以持枪之手背、撀左手手心、亦名曰𭣇。此种𭣇法、皆用三次。

\quad 繫
\quad 
疑即戏中对枪时、彼此用枪头一绕、两绕、三绕。此种打法、含有彼此枪头、繫缠在一处、不得离开的意思。按这种情形、与繫字或有相当的𨶹係。

\quad 搠
\quad 
即戏中扻或砍的意思。

\quad 籍
\quad 
即戏中揪的意思、如甲揪乙之领、乙即翻一筋斗、即为捽字之原義。

\quad 打鸳鸯场
\quad 

此种舞谱之名词、以皆在鸳鸯身着的意、戏中似无此种身段、不必强为拉扯。

按以上所说的情形、惟免有望文生训的地方、读者必定笑鄙人强为拉扯附会。然而一个人只有两手两足、古今動作、当亦无大差别、如此解法、虽然不敢说一定对、但是相去也不到其遠。由此更可能證明有机会在戏中之身段、實来源古时之舞、则毫无疑义矣。

%  第三章 論戲劇之身段

%    第一節【节】  袖谱

%    第二節【节】  手谱

%    第三節【节】  足谱

%    第四節【节】  腿谱

%    第五節【节】  腰谱

%  第四章論戲剧之身段

%    第一節【节】  鬍鬚谱

%    第二節【节】  翎子谱
\end{document}